\documentclass{article}
\usepackage[T1]{fontenc}
\usepackage{lmodern}
\usepackage{graphicx}
\usepackage{amsmath,amssymb}
\usepackage[a4paper,margin=2cm]{geometry}
\usepackage[absolute,overlay]{textpos}
\setlength{\TPHorizModule}{1mm}
\setlength{\TPVertModule}{1mm}
\usepackage[indonesian]{babel}
\usepackage{hyperref}
\setlength{\emergencystretch}{2em}
% unit: Halaman 21

\begin{document}

% === Halaman 21 (llm) ===

\begin{itemize}
  \item Hasil perhitungan frekuensi kemunculan huruf, bigram, dan trigram:
  \begin{itemize}
    \item huruf I paling sering muncul,
    \item XL adalah bigram yang paling sering muncul,
    \item XLI adalah trigram yang paling sering muncul.
  \end{itemize}
\end{itemize}

Ketiga data terbanyak ini menghasilkan dugaan bahwa

\begin{align*}
  \text{I} & \text{ berkoresponden dengan huruf plainteks e,}\
  \text{XLI} & \text{ berkoresponden dengan the,}\
  \text{XL} & \text{ berkoresponden dengan th}
\end{align*}

Pemetaan:

\begin{align*}
  \text{I} & \rightarrow \text{e}\
  \text{X} & \rightarrow \text{t}\
  \text{L} & \rightarrow \text{h}
\end{align*}

\end{document}
